\documentclass[12pt]{article}
\usepackage[utf8]{inputenc}
\usepackage[T1]{fontenc}
\usepackage[english]{babel}
\usepackage{graphicx}
\usepackage{amsmath}
\usepackage{fancyhdr}
\usepackage{geometry}
\usepackage{hyperref}
\usepackage{booktabs}
\usepackage{array}
\usepackage{float}
\usepackage{placeins}

\linespread{1.2}
\setlength{\parindent}{0.8cm}
\setlength{\parskip}{0em}
\setlength{\textfloatsep}{8pt plus 2pt minus 2pt}
\setlength{\floatsep}{8pt plus 2pt minus 2pt}
\setlength{\intextsep}{8pt plus 2pt minus 2pt}
\geometry{a4paper, portrait, margin=1in}
\setlength{\headheight}{14.5pt}
\renewcommand{\headrulewidth}{0pt}
\graphicspath{{./}}
\renewcommand{\topfraction}{0.9}
\renewcommand{\bottomfraction}{0.8}
\renewcommand{\textfraction}{0.07}
\renewcommand{\floatpagefraction}{0.8}

\newcommand{\reportfigure}[3][0.88\linewidth]{%
  \begin{figure}[!htbp]
  \centering
  \includegraphics[width=#1,height=0.30\textheight,keepaspectratio]{#2}
  \caption{#3}
  \end{figure}
}

\newcommand{\leftrotatedfigure}[2]{%
  \begin{figure}[!htbp]
  \centering
  \includegraphics[angle=90,height=0.46\textheight,keepaspectratio]{#1}
  \caption{#2}
  \end{figure}
}

\begin{document}

\begin{titlepage}
  \begin{center}
    \textsc{\large Ministry of Education of Republic of Moldova}\\[0.5cm]
    \textsc{\large Technical University of Moldova}\\[0.5cm]
    \textsc{\large Faculty of Computers, Informatics and Microelectronics}\\[0.5cm]
    \textsc{\large Software Engineering Department}\\[1.2cm]

    \vspace{20mm}
    \textsc{\Large Embedded Systems}\\[0.5cm]
    \textsc{\large Laboratory Work \#10 - Part 2}\\[0.5cm]

    \newcommand{\HRule}{\rule{\linewidth}{0.5mm}}
    \vspace{8mm}
    \HRule\\[0.4cm]
    {\LARGE \bfseries Automatic Temperature Control with Assisted Fan Cooling\\[0.3cm]}
    \HRule\\[1.5cm]

    \begin{minipage}[t]{0.45\textwidth}
      \begin{flushleft}\large
        \emph{Author:}\\
        Sava \textsc{Luchian}\\
        std. gr. FAF-233
      \end{flushleft}
    \end{minipage}
    ~
    \begin{minipage}[t]{0.45\textwidth}
      \begin{flushright}\large
        \emph{Verified:}\\
        Alexei \textsc{Martiniuc}
      \end{flushright}
    \end{minipage}

    \vspace*{\fill}
    \large Chisinau 2026
  \end{center}
\end{titlepage}

\setcounter{page}{2}
\pagestyle{fancy}
\fancyhf{}
\rhead{\thepage}
\lhead{FAF-233 Sava Luchian; Laboratory Work No.\ 10, Part 2}

\section*{Technical Task}

\subsection*{Purpose of the Laboratory}
\hspace{0.8cm} The purpose of this laboratory is to extend the temperature-control application from Part 1 with an active cooling stage. The base system remains an \textbf{Arduino Uno} application with a \textbf{10k NTC thermistor}, potentiometer setpoint, heater resistor, LCD, and \texttt{printf()} telemetry. Part 2 adds a \textbf{DC motor fan driven through L293D}; the fan speed increases gradually when the controlled resistor assembly becomes too hot and decreases gradually when the temperature returns to a safe region.

\subsection*{Objectives}
\begin{enumerate}
    \item Keep the modular ON-OFF heater controller from Part 1.
    \item Add an L293D-based fan actuator controlled with PWM.
    \item Implement gradual fan speed increase when the measured temperature exceeds the cooling threshold.
    \item Keep sequential periodic tasks rather than FreeRTOS.
    \item Display and report the final runtime values: setpoint, temperature, heater state, and fan speed.
\end{enumerate}

\subsection*{Implemented Variant}
\hspace{0.8cm} The implemented Part 2 system is based on the final Lab 9 firmware and hardware:
\begin{itemize}
    \item \textbf{Sensor input}: thermistor divider on A0.
    \item \textbf{Setpoint input}: potentiometer on A1 mapped to 20.0--80.0\,$^{\circ}$C.
    \item \textbf{Heating controller}: ON-OFF with hysteresis half-band of 1.0\,$^{\circ}$C.
    \item \textbf{Heating actuator}: resistor switched from D8.
    \item \textbf{Cooling actuator}: fan motor controlled by L293D with PWM on D11.
    \item \textbf{Reporting}: LCD + \texttt{printf()} CSV in the format \texttt{setpoint,temperature,heater,fan}.
\end{itemize}

\section{Domain Analysis}

\subsection{Technological Context and Application Domain}
\hspace{0.8cm} The thermal process in this laboratory consists of a resistive heater and nearby thermistor. Part 1 controlled the heater with ON-OFF hysteresis. Part 2 adds active cooling, which improves the response when the temperature exceeds the safe or desired range. This is a common embedded-control pattern: a simple binary heating loop is combined with proportional-like cooling effort.

\hspace{0.8cm} The L293D driver is used because the Arduino pin cannot directly drive a DC motor safely. The Uno provides direction and PWM commands, while the motor driver handles the fan current. This also keeps the motor electrically separated from the microcontroller output pins.

\subsection{Hardware Components and Justification}

\begin{table}[H]
\centering
\begin{tabular}{p{4.4cm} p{1.3cm} p{7.2cm}}
\toprule
\textbf{Component} & \textbf{Qty} & \textbf{Role in the application} \\
\midrule
Arduino Uno (ATmega328P) & 1 & Main MCU for acquisition, control, actuation, LCD handling, and serial reporting. \\
NTC thermistor 10k & 1 & Analog temperature sensing element. \\
Fixed resistor 10k & 1 & Forms the thermistor voltage divider on A0. \\
Potentiometer 10k & 1 & User-adjustable setpoint source on A1. \\
Heater resistor 220\,$\Omega$ & 1 & Resistive thermal source controlled from D8. \\
L293D motor driver & 1 & Drives the cooling fan from Uno digital pins. \\
DC motor with fan & 1 & Active cooling output with gradual speed increase. \\
LCD1602 + support resistors & 1 & Local status display for setpoint, temperature, heater, and fan state. \\
USB serial terminal / plotter & 1 & Receives \texttt{printf()} telemetry. \\
\bottomrule
\end{tabular}
\caption{Hardware components used in Lab 10 Part 2}
\label{tab:hardware}
\end{table}

\subsection{Software Components and Technologies}

\begin{itemize}
    \item \textbf{PlatformIO + VS Code}: project build and dependency management.
    \item \textbf{Arduino framework}: GPIO, ADC, timing, PWM, serial, LCD support.
    \item \textbf{AVR-libc STDIO}: \texttt{stdout} redirected to serial through \texttt{fdev\_setup\_stream}.
    \item \textbf{LiquidCrystal}: 1602 LCD driver in 4-bit mode.
    \item \textbf{Wokwi}: simulation environment for the Uno wiring and runtime checks.
\end{itemize}

\section{Conceptualisation and Design}

\subsection{Technical Requirements}

\begin{table}[H]
\centering
\small
\begin{tabular}{p{1.5cm} p{2.6cm} p{7.1cm} p{1.3cm}}
\toprule
\textbf{ID} & \textbf{Category} & \textbf{Requirement} & \textbf{Type} \\
\midrule
R-F-01 & Sensor & System shall read temperature from a thermistor divider. & Func. \\
R-F-02 & Setpoint & System shall read setpoint from potentiometer on A1. & Func. \\
R-F-03 & Heating & System shall implement ON-OFF hysteresis for the heater. & Func. \\
R-F-04 & Cooling & System shall drive a fan through L293D when temperature is too high. & Func. \\
R-F-05 & Cooling & Fan speed shall increase gradually as temperature rises above the cooling threshold. & Func. \\
R-F-06 & Display & LCD shall show setpoint, temperature, heater state, and fan level. & Func. \\
R-F-07 & Reporting & Runtime values shall be sent through \texttt{printf()} as CSV. & Func. \\
R-NF-01 & Timing & The application shall use sequential periodic tasks. & Non-func. \\
R-NF-02 & Modularity & Code shall be split into reusable modules. & Non-func. \\
R-NF-03 & Platform & Firmware shall fit Arduino Uno resource limits. & Non-func. \\
\bottomrule
\end{tabular}
\caption{Technical requirements for Lab 10 Part 2}
\label{tab:requirements}
\end{table}

\subsection{Control Design}

\hspace{0.8cm} The heater is controlled by hysteresis:
\[
Heater =
\begin{cases}
ON, & T \leq SP - H \\
OFF, & T \geq SP + H \\
previous, & SP - H < T < SP + H
\end{cases}
\]

\hspace{0.8cm} The fan is controlled by a gradual cooling law. It remains off below the cooling start threshold, then increases toward full speed as the measured temperature rises:
\[
Fan_{target} =
\begin{cases}
0\%, & T \leq SP + 1.0 \\
100 \cdot \frac{T-(SP+1.0)}{9.0}, & SP+1.0 < T < SP+10.0 \\
100\%, & T \geq SP + 10.0
\end{cases}
\]

\hspace{0.8cm} The actual fan output does not jump immediately to the target. It changes by 5\% every 100\,ms, so the fan accelerates and decelerates gradually.

\subsection{Timing Design}

\begin{table}[H]
\centering
\begin{tabular}{llll}
\toprule
\textbf{Activity} & \textbf{Period} & \textbf{Trigger} & \textbf{Handler} \\
\midrule
Setpoint update & 100\,ms & \texttt{millis()} elapsed & \texttt{runSetpointTask()} \\
Acquisition update & 100\,ms & \texttt{millis()} elapsed & \texttt{runAcquisitionTask()} \\
Control update & 250\,ms & \texttt{millis()} elapsed & \texttt{runControlTask()} \\
Report update & 500\,ms & \texttt{millis()} elapsed & \texttt{runReportTask()} \\
Fan ramp update & 100\,ms & \texttt{millis()} elapsed & fan actuator \texttt{tick()} \\
\bottomrule
\end{tabular}
\caption{Sequential periodic timing model for Lab 10 Part 2}
\label{tab:timing}
\end{table}

\section{Hardware and Software Implementation}

\subsection{Electrical Schematic and Hardware Setup}
\reportfigure[0.94\linewidth]{electrical_schema_10.png}{Electrical circuit schematic for Lab 10 Part 2}
\leftrotatedfigure{hardware1.jpg}{Hardware setup while the heater resistor is active}
\leftrotatedfigure{hardware2.jpg}{Hardware setup after the heater resistor is switched off}

\textbf{Pin assignment:}
\begin{itemize}
    \item A0 -- thermistor divider node
    \item A1 -- potentiometer signal
    \item D8 -- heater resistor output
    \item D2 -- L293D IN1
    \item D10 -- L293D IN2
    \item D11 -- L293D EN1,2 (PWM)
    \item D7 -- LCD RS
    \item D6 -- LCD E
    \item D5 -- LCD D4
    \item D4 -- LCD D5
    \item D3 -- LCD D6
    \item D13 -- LCD D7
    \item D0/D1 -- USB serial UART for STDIO output
\end{itemize}

\subsection{Project Directory Structure}

\begin{verbatim}
lab10/
|-- report/
|   |-- main.tex
|   |-- electrical_schema_10.png
|   |-- hardware1.jpg
|   |-- hardware2.jpg
|   `-- terminal.png
\end{verbatim}

\hspace{0.8cm} The firmware used for Part 2 is the final Lab 9 implementation after adding the L293D fan actuator and gradual fan-speed control.

\subsection{Software Module Descriptions}

\begin{itemize}
    \item \textbf{ThermistorSensor}: converts ADC readings to voltage, resistance, and temperature.
    \item \textbf{PotentiometerInput}: maps ADC value to the configurable setpoint range.
    \item \textbf{OnOffHysteresisController}: implements thresholded heater logic with memory.
    \item \textbf{HeaterActuator}: drives the heating resistor from a digital pin.
    \item \textbf{L293dFanActuator}: drives and ramps the cooling fan through L293D using PWM.
    \item \textbf{CommandInterface}: handles serial commands such as \texttt{help} and \texttt{status}.
    \item \textbf{LcdDisplay}: presents local runtime status on the 1602 LCD.
    \item \textbf{Lab9ControlApp}: orchestrates sequential periodic tasks used by Part 2.
\end{itemize}

\section{Testing and Validation}

\subsection{Build Results}

\begin{table}[H]
\centering
\begin{tabular}{lrrrr}
\toprule
\textbf{Environment} & \textbf{Flash used} & \textbf{Flash \%} & \textbf{RAM used} & \textbf{RAM \%} \\
\midrule
\texttt{uno} & 11\,430 B & 35.4\% & 788 B & 38.5\% \\
\bottomrule
\end{tabular}
\caption{PlatformIO build resource usage for the Part 2 firmware}
\label{tab:build}
\end{table}

\subsection{Terminal Output}
\reportfigure[0.90\linewidth]{terminal.png}{Serial terminal output for the Part 2 implementation}

\subsection{Test Scenarios and Execution Results}

\begin{table}[H]
\centering
\small
\begin{tabular}{p{0.8cm} p{3.2cm} p{6.0cm} p{1.3cm}}
\toprule
\textbf{ID} & \textbf{Scenario} & \textbf{Observed result} & \textbf{Result} \\
\midrule
T-01 & Startup and build & Firmware builds and starts correctly on Uno profile. & PASS \\
T-02 & Setpoint sweep & Potentiometer changes setpoint continuously across the configured range. & PASS \\
T-03 & Heater low temperature case & When \(T \le SP-H\), the heater turns ON. & PASS \\
T-04 & Heater high temperature case & When \(T \ge SP+H\), the heater turns OFF. & PASS \\
T-05 & Fan start threshold & Above \(SP+1.0\,^{\circ}C\), the fan begins to spin. & PASS \\
T-06 & Fan ramp behaviour & Fan speed increases progressively as temperature rises. & PASS \\
T-07 & Fan slowdown behaviour & Fan speed decreases progressively as the temperature falls. & PASS \\
T-08 & LCD/serial output & LCD and CSV output reflect current setpoint, temperature, heater, and fan state. & PASS \\
T-09 & Resource limits & Flash and RAM remain below Uno limits. & PASS \\
\bottomrule
\end{tabular}
\caption{Test execution summary for Lab 10 Part 2}
\label{tab:testresults}
\end{table}

\subsection{Representative Serial Output}

\begin{verbatim}
 35.0, 24.8,1,0
 35.0, 29.2,0,0
 35.0, 38.4,0,26
 35.0, 42.2,0,67
 35.0, 45.1,0,100
\end{verbatim}

\hspace{0.8cm} The four channels represent setpoint, measured temperature, heater state, and fan speed percentage.

\section*{General Conclusions}
\hspace{0.8cm} Laboratory Work \#10 Part 2 extended the Part 1 thermistor-heater controller with active fan cooling. The L293D driver enabled safe motor control, while PWM ramping avoided abrupt fan-speed changes. The resulting system remains sequential, modular, observable through LCD and serial telemetry, and small enough for Arduino Uno.

Key outcomes:
\begin{enumerate}
    \item The heater still follows ON-OFF hysteresis around the potentiometer-controlled setpoint.
    \item The fan provides an additional cooling channel above the high-temperature region.
    \item The fan speed changes gradually, which reduces mechanical and electrical abruptness.
    \item The full system remains compatible with PlatformIO, Wokwi, and Arduino Uno limits.
\end{enumerate}

\section*{Note on AI Tool Usage}
\hspace{0.8cm} AI-assisted tools were used to support coding productivity and report drafting:
\begin{itemize}
    \item \textbf{GitHub Copilot} -- code completion and refactoring assistance.
    \item \textbf{LLM assistance} -- LaTeX drafting, editing support, and technical structure refinement.
\end{itemize}
All generated content was manually reviewed and adapted to the implemented project.

\section*{Bibliography}
\begin{enumerate}
    \item Arduino Official Reference. \url{https://www.arduino.cc/reference/en/}
    \item AVR-libc User Manual. \url{https://www.nongnu.org/avr-libc/user-manual/}
    \item PlatformIO Documentation. \url{https://docs.platformio.org/en/latest/}
    \item Wokwi Simulator Documentation. \url{https://docs.wokwi.com/}
    \item LiquidCrystal Library Reference. \url{https://www.arduino.cc/reference/en/libraries/liquidcrystal/}
    \item L293D Motor Driver Datasheet. Texas Instruments / STMicroelectronics.
\end{enumerate}

\FloatBarrier
\section*{Annex -- Source Code and Report Assets}

The final Part 2 firmware is based on the Lab 9 implementation with the motor fan extension.

\bigskip
\noindent\textbf{Repository URL:}\\
\url{https://github.com/Ekkusuu/embedded-systems-repo/tree/main/lab9}

\bigskip
\noindent Lab 10 report assets are stored in:\\
\texttt{lab10/report/}

\end{document}
